\documentclass[conference]{IEEEtran}

\usepackage{cite}
\usepackage{amsmath,amssymb,amsfonts}
\usepackage{graphicx}
\usepackage{booktabs}
\usepackage{multirow}

\begin{document}

\title{IoT-Based Autonomous Mobile Robot for Path Planning and Real-Time Obstacle Avoidance Using ESP32}

\author{
\IEEEauthorblockN{Ankit Raj, Reesha C. Koli, Sushil S. Chandaragi, Shruti Jadhav, Sejal Halgekar}
\IEEEauthorblockA{
Department of Computer Science and Engineering\\
KLE Technological University's Dr. M. S. Sheshgiri College of Engineering and Technology\\
Belagavi, Karnataka, India\\
}
}

\maketitle

\begin{abstract}
This paper presents the design and implementation of a low-cost, IoT-enabled autonomous mobile robot for indoor transportation. The system achieves PID-based line following, junction-aware node navigation, and real-time obstacle avoidance. Built around an ESP32 microcontroller, it utilizes a five-channel infrared sensor array for path tracking, an HC-SR04 ultrasonic sensor mounted on an SG90 servo for directional obstacle scanning, an MPU6050 IMU for stability, and an L298N motor driver. The robot communicates via the MQTT protocol and is managed through a real-time React-based web dashboard. Experimental evaluation demonstrates reliable navigation with line-following accuracy of $\sim$91\%, obstacle detection accuracy of $\sim$90\%, and average obstacle response time under 600~ms. The proposed architecture provides a scalable, affordable solution for autonomous indoor logistics.
\end{abstract}

\begin{IEEEkeywords}
Autonomous Mobile Robot, IoT, ESP32, PID Control, MQTT, Obstacle Avoidance, Indoor Navigation.
\end{IEEEkeywords}

\section{Introduction \& Related Work}
Autonomous mobile robots are increasingly deployed in hospitals, warehouses, and smart buildings to automate repetitive indoor transportation tasks. A major challenge in this domain is accurately following predefined paths while avoiding obstacles using resource-constrained hardware. Advanced systems based on computer vision, SLAM, and deep learning \cite{ref2,ref6,ref8} show promising results but demand high computational resources and expensive hardware. Conversely, traditional approaches like the Artificial Potential Field \cite{ref9} and Dynamic Window Approach \cite{ref5} are computationally lighter but often lack remote monitoring capabilities.

To bridge this gap, this work proposes a low-cost autonomous indoor courier robot that integrates PID-based line following, sensor fusion, servo-mounted ultrasonic obstacle scanning, and real-time IoT monitoring. The system uses commercially available off-the-shelf components, offering stable navigation, real-time obstacle avoidance, junction-aware delivery routing, and a full MQTT-based IoT dashboard for remote tuning and control. Compared to existing solutions \cite{ref1,ref4,ref7}, our approach provides functional performance and rich remote operability at a fraction of the cost ($\sim$\$30 USD).

\section{Proposed System Architecture}
The system architecture (Fig.~\ref{fig:architecture}) is divided into the edge hardware, an MQTT broker, and a web dashboard. 

\begin{figure}[ht]
\centering
\includegraphics[width=0.48\textwidth]{iot methodology.jpeg}
\caption{Proposed System Architecture}
\label{fig:architecture}
\end{figure}

\subsection{Hardware Configuration \& Power Management}
The robot (Fig.~\ref{fig:prototype}) is built on a 4WD chassis driven by four DC motors via an L298N dual H-bridge motor driver. An ESP32 serves as the central controller, gathering data from a five-channel IR sensor array (mounted underneath), an MPU6050 IMU, and an HC-SR04 ultrasonic sensor mounted on an SG90 servo for 150\textdegree{} directional scanning. Power is supplied by a 2S (7.4V) lithium-ion battery pack with a TP4056 charger and 2S BMS. To prevent brownout resets during motor current spikes, the battery powers the L298N driver directly, while an LM2596 buck converter steps the voltage down to 5V to power the ESP32 and logic sensors.

\begin{figure}[ht]
\centering
\includegraphics[width=0.45\textwidth]{model.jpeg}
\caption{Developed Autonomous Mobile Robot Prototype}
\label{fig:prototype}
\end{figure}

\subsection{IoT Communication and Web Dashboard}
The system employs an MQTT architecture with a Node.js (Aedes) broker. The ESP32 publishes sensor data, motor telemetry, obstacle distances, and radar sweeps at 2--5~Hz. A React-based web dashboard connects via WebSockets, providing real-time visualization of IR sensor states, motor PWMs, obstacle radar, and navigation status. The dashboard allows operators to tune PID parameters ($K_p, K_i, K_d$), adjust base speed, manually control the robot, and issue destination routes (e.g., ``Ward 1'').

\section{Navigation Methodology}
The robot's autonomous operation is governed by a finite state machine integrating line following, junction handling, and obstacle avoidance.

\subsection{PID-Based Line Following \& Junction Navigation}
The five-channel IR sensor array provides binary line detection. Sensors are assigned positional weights $w_i \in \{-2, -1, 0, 1, 2\}$. The navigation error $e(t)$ is the weighted average of active sensors. A PID controller computes the steering correction $u(t) = K_p e(t) + K_i \int e(t) dt + K_d \frac{de(t)}{dt}$. Left and right motor speeds are adjusted by adding/subtracting $u(t)$ from a base speed. 

A T-junction is detected when both outer-left and outer-right sensors simultaneously detect the line. Upon detection, the robot pops the next command (e.g., Turn Left, Turn Right, Stop) from a predefined route queue. Turns are executed as timed pivot maneuvers (700~ms), aided by MPU6050 directional feedback.

\subsection{Obstacle Detection and Avoidance}
The forward-facing ultrasonic sensor measures distance at 2~Hz to minimize interference with the PID loop. To prevent false positives, three consecutive readings below 15~cm are required to trigger an obstacle state. Upon confirmation:
1) The robot stops and performs a radar sweep, rotating the servo from 15\textdegree{} to 165\textdegree{} in 5\textdegree{} increments.
2) Sweep data is published as a batch MQTT message for the dashboard's radar visualizer.
3) The robot compares left ($D_L$) and right ($D_R$) clearances.
4) It executes a pivot turn toward the direction of greater clearance.
5) It drives forward and searches for the original line using the IR sensors, resuming PID navigation upon detection.

\section{Experimental Results}
The robot was tested in an indoor environment on black tape tracks featuring straight segments, curves, T-junctions, and static obstacles. 

\begin{table}[ht]
\caption{Overall System Performance Metrics}
\label{tab:performance}
\centering
\begin{tabular}{@{}llc@{}}
\toprule
\textbf{Category} & \textbf{Metric} & \textbf{Success/Value} \\ \midrule
\multirow{3}{*}{\textbf{Path Tracking}} 
& Straight Path Accuracy & 95\% (20 trials) \\
& Curved Path Accuracy & 93\% (15 trials) \\
& T-Junction Navigation & 90\% (20 trials) \\ \midrule
\multirow{4}{*}{\textbf{Obstacle Handling}} 
& Detection Accuracy & 90\% (10 trials) \\
& Avoidance Success Rate & 80\% (10 trials) \\
& Average Response Time & $\sim$600 ms \\
& Radar Sweep Duration & $\sim$2.5 s \\ \midrule
\textbf{System} 
& Dashboard Latency & $<$100 ms \\
& Battery Duration & 2--3 hours \\ \bottomrule
\end{tabular}
\end{table}

As shown in Table~\ref{tab:performance}, the system achieved robust line tracking (91\% overall accuracy). Performance slightly decreased on sharp 90\textdegree{} curves due to the open-loop nature of timed pivot turns. The obstacle detection algorithm effectively suppressed noise, achieving 90\% accuracy with a response time of $\sim$600~ms. The avoidance maneuver succeeded in 80\% of trials, with occasional failures occurring during the line re-acquisition phase. The real-time dashboard proved invaluable for rapidly tuning PID gains without reflashing firmware.

\section{Conclusion}
This paper presented an affordable, IoT-enabled autonomous mobile robot capable of PID line following, junction routing, and servo-assisted obstacle avoidance. Utilizing an ESP32 and an MQTT-based web dashboard, the system provides reliable navigation and remote teleoperation at a low hardware cost. Future work will focus on integrating camera-based lane detection, SLAM for map-free navigation, and multi-robot coordination for advanced smart logistics applications.

\section*{Acknowledgment}
The authors thank Dr. Niranjan Muchandi and the Dept. of CSE, Dr. M. S. Sheshgiri College of Engineering and Technology, KLE Tech. University, Belagavi, for their guidance.

\begin{thebibliography}{00}
\bibitem{ref1} P. Kotikalapudi et al., ``Obstacle Avoidance and Path Finding for Mobile Robot Navigation,'' \textit{CS \& IT}, 2020.
\bibitem{ref2} J. Hong et al., ``Semantically-Aware Strategies for Stereo-Visual Robotic Obstacle Avoidance,'' \textit{ICRA}, 2021.
\bibitem{ref4} Y. Chen et al., ``SLP-Improved DDPG Path Planning Algorithm for Mobile Robots,'' \textit{Sensors}, 2023.
\bibitem{ref5} Y. Sun et al., ``Local Path Planning for Mobile Robots Based on Fuzzy Dynamic Window Algorithm,'' \textit{Sensors}, 2023.
\bibitem{ref6} J. Gao et al., ``Deep Reinforcement Learning for Indoor Mobile Robot Path Planning,'' \textit{Sensors}, 2020.
\bibitem{ref7} K. Almazrouei et al., ``Dynamic Obstacle Avoidance and Path Planning through Reinforcement Learning,'' \textit{Appl. Sci.}, 2023.
\bibitem{ref8} Y. Zhang et al., ``Path Planning of a Mobile Robot for Dynamic Indoor Environment Based on SAC-LSTM,'' \textit{Sensors}, 2023.
\bibitem{ref9} O. Khatib, ``Real-Time Obstacle Avoidance for Manipulators and Mobile Robots,'' \textit{IJRR}, 1986.
\end{thebibliography}
\end{document}
